
\section{问题重述}

\subsection{问题背景}

% 第1段：宏观背景（2-3句话，从社会发展趋势/技术需求引入）
% 示例：随着全球能源问题和环境问题日益严峻，人们迫切需要开发一种新能源来替代传统的化石燃料，以此来达到"碳达峰""碳中和"的目标。XXX作为一种分布最广、来源最稳定以及储量最大的新能源受到了人们的密切关注。

% 第2段：问题来源/原因（2-3句话，说明研究对象及其重要性）
% 示例：XXX是一种较为理想的、技术发展相对成熟的大规模利用XXX技术，XXX是其XXX的重要基本组件。

% 第3段：引出待解决问题（1句话）
% 示例：如何通过数学建模对XXX的各项参数进行优化设计使得XXX最大成为了我们亟待解决的问题。

\subsection{问题重述}

% 第1段：系统描述（2-3句话，描述研究对象的组成/工作原理）
% 示例：XXX的底座包括XXX和XXX，其中XXX可以控制XXX，XXX可以控制XXX。随着XXX的改变，XXX也跟着改变，使得XXX可以XXX。XXX内排列了大量的XXX以及一个XXX，XXX位于XXX，其以XXX的形式储存吸收得到的XXX，然后再进一步将XXX转换成XXX。

% 逐问重述（每问一段，不加粗"问题X："）
\textbf{问题一：}在题目给定的【范围/条件】内，【已知条件1】，【已知条件2】，且附件中给出了【数据】，计算该【研究对象】产生的【指标1】、【指标2】以及【指标3】。

\textbf{问题二：}根据题目给定的要求，【约束条件1】，在【统一条件】的条件下，给出该【研究对象】中【决策变量1】的位置以及【决策变量2】的位置、数目、【参数1】和【参数2】，使得该【研究对象】在满足【约束】的同时，能够使【目标】尽量大。

\textbf{问题三：}【约束条件】仍为【数值】，但此时【变化描述】均可以不同，在此条件下，重新给出【研究对象】中【决策变量1】的位置以及【决策变量2】的位置、【决策变量3】的总数目、各个【决策变量4】的【参数1】以及【参数2】，使得该【研究对象】在满足【约束】的同时，能使【目标】尽量大。
